% ECONOMICS_PROBLEM: {"id": "D83-480", "title": "Unknown unknowns", "jel_code": "D83", "source_id": "OP-0276"}
% !TeX program = xelatex
\documentclass[11pt,a4paper]{article}
\usepackage[margin=23mm]{geometry}
\usepackage{fontspec}
\setmainfont{Latin Modern Roman}
\usepackage{amsmath,amssymb,mathtools}
\usepackage{xcolor,enumitem,booktabs,longtable,array,xurl,hyperref}
\definecolor{muted}{HTML}{53636C}
\hypersetup{colorlinks=true,urlcolor=blue,linkcolor=blue}
\setlength{\parindent}{0pt}
\setlength{\parskip}{5pt}
\newcommand{\R}{\mathbb R}
\newcommand{\E}{\mathbb E}
\newcommand{\N}{\mathbb N}
\newcommand{\ind}{\mathbf 1}
\newcommand{\Var}{\operatorname{Var}}
\newcommand{\Cov}{\operatorname{Cov}}
\newcommand{\supp}{\operatorname{supp}}
\DeclareMathOperator*{\argmin}{arg\,min}
\DeclareMathOperator*{\argmax}{arg\,max}
\newcommand{\entrymeta}[3]{{\sffamily\small\color{muted}\textbf{\texttt{#1}}\quad|\quad #2\par #3\par}}
\newcommand{\Problemref}[1]{\texttt{#1}}
\newcommand{\JELref}[2]{\texttt{#2} (JEL \texttt{#1})}
\begin{document}
\section*{Unknown unknowns}
\textbf{Problem ID:} \texttt{D83-480}\quad \textbf{JEL:} \texttt{D83}\par
% BEGIN ECONOMICS STATEMENT
\label{problem:OP-0276}
{\sffamily\small\color{muted}Original subject group: Economic theory, equilibrium and social choice\par}
\label{definitions:problem:30007557}
\textbf{Audit classification:} Background; proposed extension. The genuine October 2026 paper provides representation results. Its subjective unknowns are not the finite analyst-partition experiment proposed here; no exact remaining identification question was verified.
\subsection*{Definitions and precise research target}
\textbf{Model and research target.} Let \(S\) be a finite latent state space used by the analyst, \(X\) a finite consequence set, and \(\mathcal P_0\preceq\mathcal P_1\) two partitions of \(S\), with the second refining the first. The decision maker initially describes events only through \(\mathcal P_0\) and knows that additional distinctions may later become expressible. Observe choices from menus of acts \(a:S\to\Delta(X)\), first through coarse descriptions and then after randomized disclosure of the refinement. Here \(\Delta(X)\) denotes lotteries over consequences. Let \(\mathcal R\) be a stated class of representations distinguishing subjective likelihoods, consequence utility, and attitudes toward awareness expansion. For observed choice law \(P\), define
\[
 \mathcal R(P)=\{r\in\mathcal R:r\text{ generates }P\}.
\]
The proposed target is to identify which attitudes toward indescribable contingencies are recoverable from these menu choices and which remain confounded with ordinary risk or ambiguity preferences. Derive sharp identified sets and experimental restrictions that separate the components. For a policy offering act \(a\) rather than \(b\), report welfare comparisons robust across \(\mathcal R(P)\), conditional on each representation's stated welfare interpretation. Analyze sensitivity to the analyst's chosen latent refinement.

\emph{Definitions and maintained assumptions (editorial unless a source definition is named).} For partitions, write \(\mathcal P_0\preceq\mathcal P_1\) when every cell of \(\mathcal P_1\) lies in a cell of \(\mathcal P_0\). An initially describable act must be \(\mathcal P_0\)-measurable, hence constant on its cells; it is incoherent to present a genuinely fine-state-contingent contract using only its coarse description without specifying how unresolved cells pay. One possible editorial experiment presents coarse measurable acts initially and newly available \(\mathcal P_1\)-measurable acts after disclosure. Utility is a function \(u:X\to\mathbb R\), lotteries belong to \(\Delta(X)\), and each representation specifies probabilities over expressible events plus a rule for assessing future refinements. An analyst’s finite latent \(S\) is not a claim that the participant initially conceptualizes all its states. A surprise disclosure is informative as well as awareness-expanding; to separate these channels, specify matched information-only controls. Observational fit is equality of choice laws under every recorded menu and disclosure condition; robust welfare requires a nonempty representation class and a stated normative interpretation.
\subsection*{Variants and assumption changes}
\begin{enumerate}
\item \textbf{Known expansion (editorial).} Tell participants that a particular coarse cell will be split, withholding the realized finer state. Identify preference for access to contingent acts separately from utility and probability parameters.
\item \textbf{Unanticipated expansion (editorial).} Do not announce the new contingencies, and compare post-disclosure behavior with matched information-only controls. Determine what, if anything, is identified about pre-disclosure unawareness; finite-state evidence cannot certify attitudes to all genuinely indescribable contingencies.
\end{enumerate}
\subsection*{References and source audit}
\begin{enumerate}
\item Full publisher PDF and October1,2026 date verified; representation results only. Exact proposed experiment absent from checked support. \url{https://link.springer.com/article/10.1007/s00199-026-01750-z}
\end{enumerate}
\begin{enumerate}\raggedright
\item\hypertarget{definitions:source:30007557:1}{} Marie-Louise Vier\o{}. 2026. \emph{Lost in objective translation: awareness of unawareness when unknowns are not simply unknowns}. Abstract and representation results.
\url{https://link.springer.com/article/10.1007/s00199-026-01750-z}.
DOI: \url{https://doi.org/10.1007/s00199-026-01750-z}.
\end{enumerate}

\subsection*{Common conventions: Source mathematical conventions}

These conventions apply to each entry unless explicitly replaced.

\textbf{Models and identification.} A primitive model \(m\) specifies all probability laws, feasible choices, information, equilibrium and measurement rules used by its target. The admissible class \(\mathcal M\) is fixed independently of outcomes. If \(P_O(m)\) is its observable probability law and \(\tau(m)\) the target, its identified set at observable law \(P\) is
\[
 \mathcal I_\tau(P)=\{\tau(m):m\in\mathcal M,\ P_O(m)=P\}.
\]
Point identification means that this set is a singleton. An empty set signals incompatibility of data and assumptions. Coordinate bounds need not describe the joint set. Bounds are sharp only if every retained target is attainable or arbitrarily approachable by compatible models. Finite bounds require explicit restrictions on unobserved quantities; unrestricted confounding, hidden wealth, extrapolation or arbitrary equilibrium selection can yield uninformative sets.

\textbf{Causal interventions.} A potential outcome \(Y_i(p)\) is the outcome under a complete policy regime \(p\), including timing, financing, information and market spillovers. The target population and units are fixed before treatment. With interference, use the complete assignment vector or a justified exposure mapping; individual no-interference notation alone cannot identify an economy-wide intervention. A difference of mean potential outcomes does not require the cross-world joint distribution. Conditioning on post-treatment migration, survival, enrollment or employment changes the estimand and needs separate justification.

\textbf{Statistical statements.} An estimator, confidence set, forecast or test specifies a sampling law, model class and loss/coverage criterion. Uniform \((1-\alpha)\)-coverage means \(\inf_{m\in\mathcal M}\Pr_m\{\tau(m)\in C_n\}\geq1-\alpha\), or an explicitly asymptotic version. The whole parameter space trivially has coverage, so informativeness requires a stated width, risk or power objective. A forecasting improvement is relative to a named benchmark, information set and evaluation population.

\textbf{Equilibrium and optimization.} An equilibrium comprises feasible plans, optimality, beliefs where needed, budget constraints and all market/resource clearing. The primitives or a stated benchmark define each correspondence \(\mathcal E(p;m)\). Nonemptiness is assumed or proved, never inferred just from compact policy sets. A selected-equilibrium target uses a declared selection rule; a robust target quantifies over all permitted equilibria. Use suprema or infima unless attainment is established. An \(\varepsilon\)-optimal policy is feasible and within \(\varepsilon\) of the finite optimum value. Report an empty feasible set separately.

\textbf{Welfare and accounting.} Normative utility, interpersonal normalization, social weights, numeraire, discounting and the use of public receipts are primitives. Behavioral utility need not coincide with the welfare criterion. Household welfare includes ownership income and rents once; adding profits again to compensating variation generally double-counts them. A monetary equivalent is defined by an explicit transfer and price convention, with monotonicity and range conditions. Average effects, mechanism contributions, transfers and welfare losses are different targets. Infinite sums require convergence and dynamic feasible plans require terminal or no-Ponzi conditions.

\textbf{Source versions.} A working-paper date, revision date and journal publication year are distinguished. A DOI for a journal version is not automatically the DOI of an earlier manuscript. Locators refer to the stated version and printed pages where specified. A metadata-only check does not verify a claim about a full-text conclusion. Every entry's source subsection narrows the provenance claim accordingly.
% END ECONOMICS STATEMENT
\end{document}
